
\section*{Глава 3. Разработка веб-приложения}
\addcontentsline{toc}{section}{Глава 3. Разработка веб-приложения}
\stepcounter{section}

\subsection{Разработка базовой структуры приложения и вёрстка шаблонов страниц}

В соответствии с задачей~3.1, была реализована клиентская часть веб-приложения
FitCenter на базе библиотеки \textbf{React.js} с использованием маршрутизатора
\textbf{React Router}. Приложение организовано по компонентному принципу:
общие элементы интерфейса вынесены в каталог \texttt{components/}, а страницы
приложения --- в каталог \texttt{pages/}. Центральный файл \texttt{App.js}
определяет маршруты и обеспечивает единую навигационную панель для всех
разделов сайта.

Для разграничения доступа к страницам реализован компонент-обёртка
\texttt{PrivateRoute}, который проверяет наличие авторизованного пользователя
и соответствие его роли списку допустимых ролей. Публичными являются страницы
расписания (\texttt{/}), входа и регистрации (\texttt{/login}, \texttt{/register}),
а также информационные разделы «Тренеры» и «О нас». Личный кабинет клиента
доступен по маршруту \texttt{/profile}, кабинет тренера --- по маршрутам
\texttt{/trainer} и \texttt{/trainer/profile}, панель администратора --- по
маршрутам \texttt{/admin/*}.

Вёрстка страниц выполнена с использованием единого файла стилей
\texttt{index.css}. Для отображения расписания применяется компонент карточки
тренировки, который выводит фотографию, название, дату, время, имя тренера и
количество свободных мест. Пример структуры карточки представлен
в листинге~\ref{lst:workout_card}.
\clearpage

\begin{lstlisting}[caption={Фрагмент компонента карточки тренировки (Schedule.js, стр.~8--41)}, label={lst:workout_card}]
function WorkoutCard({ workout, onBook, view }) {
  const { user } = useAuth();
  const active = workout.bookings?.filter(b => b.status === 'active') || [];
  const spotsLeft = workout.capacity - active.length;
  const isFull = spotsLeft <= 0;
  const bookedCount = Math.min(active.length, workout.capacity);
  const isBooked = active.some(b => String(b.user_id) === String(user?.id));
  const isPast = new Date(workout.starts_at) < new Date();

  const BookBtn = () => {
    if (!user) {
      return (
        <button
          className="btn btn-primary btn-full"
          onClick={() => { window.location.href = '/login'; }}
        >
          Войти для записи
        </button>
      );
    }
    if (user.role !== 'client') return null;
    return (
      <button className="btn btn-primary btn-full"
        disabled={isFull || isBooked || isPast}
        onClick={() => onBook(workout.id)}>
        {isBooked ? ' Вы записаны' : isFull ? 'Мест нет' : isPast ? 'Прошла' : 'Записаться'}
      </button>
    );
  };
\end{lstlisting}

Серверная часть приложения реализована на языке \textbf{Go} с фреймворком
\textbf{Gin} и организована в виде слоёв: маршруты (\texttt{transport/}),
обработчики HTTP-запросов (\texttt{handlers/}), бизнес-логика (\texttt{service/})
и модели данных (\texttt{models/}). Такое разделение упрощает сопровождение
кода и соответствует архитектуре, описанной в разделе~2.

\subsection{Реализация аутентификации, авторизации и системы ролей (клиент, тренер, администратор)}

Согласно задаче~3.2, в системе реализована ролевая модель доступа с тремя
ролями: \texttt{client} (клиент), \texttt{trainer} (тренер) и \texttt{admin}
(администратор). Для обеспечения безопасности используется протокол
\textbf{JWT} (JSON Web Token). При успешном входе или регистрации сервер
генерирует токен, содержащий идентификатор пользователя (\texttt{user\_id}) и
поле \texttt{role}, определяющее права доступа. Пример формирования токена
представлен в листинге~\ref{lst:jwt}.

\begin{lstlisting}[caption={Генерация JWT-токена при аутентификации (jwtutil.go, стр.~11--18)}, label={lst:jwt}]
func GenerateToken(user models.User) (string, error) {
	claims := jwt.MapClaims{
		"user_id": user.ID.String(),
		"role":    string(user.Role),
		"exp":     time.Now().Add(72 * time.Hour).Unix(),
	}
	token := jwt.NewWithClaims(jwt.SigningMethodHS256, claims)
	return token.SignedString([]byte(os.Getenv("JWT_SECRET")))
}
\end{lstlisting}

Пароли пользователей не хранятся в открытом виде: при регистрации и создании
учётной записи администратором применяется алгоритм хеширования
\textbf{bcrypt}. Самостоятельная регистрация в роли тренера запрещена ---
тренеры создаются только администратором, что предотвращает несанкционированное
получение расширенных прав.

На стороне сервера реализованы middleware-функции \texttt{AuthMiddleware} и
\texttt{RequireRole}. Первая проверяет наличие и корректность JWT-токена
в заголовке \texttt{Authorization}, вторая --- соответствие роли пользователя
требованиям конкретного маршрута. На клиенте контекст \texttt{AuthContext}
сохраняет токен в \texttt{localStorage}, автоматически подставляет его в
запросы к API и восстанавливает сессию при перезагрузке страницы.

Маршруты API сгруппированы по уровню доступа (листинг~\ref{lst:routes}):
публичные эндпоинты доступны без авторизации, защищённые --- после проверки
токена, а операции управления тренировками и администрирования --- только
для ролей \texttt{trainer} и \texttt{admin} соответственно.

\begin{lstlisting}[caption={Группировка маршрутов API по ролям (routes.go, стр.~30--43)}, label={lst:routes}]
	trainer := protected.Group("/")
	trainer.Use(RequireRole("trainer", "admin"))
	trainer.POST("/workouts", h.CreateWorkout)
	trainer.PUT("/workouts/:id", h.UpdateWorkout)
	trainer.DELETE("/workouts/:id", h.DeleteWorkout)
	trainer.POST("/workouts/:id/photo", h.UploadWorkoutPhoto)

	adminRoutes := protected.Group("/admin")
	adminRoutes.Use(RequireRole("admin"))
	adminRoutes.GET("/users", h.ListUsers)
	adminRoutes.POST("/users", h.CreateUser)
	adminRoutes.PUT("/users/:id", h.UpdateUser)
	adminRoutes.DELETE("/users/:id", h.DeleteUser)
	adminRoutes.GET("/stats", h.GetStats)
\end{lstlisting}

\subsection{Разработка CRUD-интерфейса для управления тренировками, расписанием и пользователями}

В соответствии с задачей~3.3, реализован полный набор операций создания,
чтения, обновления и удаления (CRUD) для ключевых сущностей системы.

Управление тренировками доступно тренеру в кабинете \texttt{TrainerDashboard}
и администратору в разделе «Тренировки» панели \texttt{AdminDashboard}.
Тренер может создавать, редактировать и удалять только собственные занятия;
администратор --- тренировки любого тренера. При создании и изменении
тренировки проверяется, что дата и время начала не относятся к прошедшему
периоду. Данные передаются на сервер через REST API (\texttt{POST /api/workouts},
\texttt{PUT /api/workouts/:id}, \texttt{DELETE /api/workouts/:id}).

Просмотр расписания реализован на странице \texttt{Schedule}. Клиент и
неавторизованный пользователь могут фильтровать тренировки по дате и тренеру.
Список формируется запросом \texttt{GET /api/workouts} с параметрами
\texttt{date} и \texttt{trainer\_id}. Сервер возвращает тренировки вместе
с данными о тренере и активных записях, что позволяет сразу отображать
заполненность занятий.

Управление пользователями реализовано в панели администратора. Администратор
просматривает список всех учётных записей, создаёт новых пользователей с
заданной ролью, изменяет имя, email и роль, а также удаляет пользователей.
При удалении тренера каскадно удаляются связанные тренировки и записи на них,
что обеспечивает целостность данных в базе.

\subsection{Реализация онлайн-записи на тренировку с проверкой доступных мест и запретом двойной записи}

Задача~3.4 реализована в сервисном слое backend-части методом \texttt{Book},
который выполняет запись клиента на выбранную тренировку. Перед созданием
записи система последовательно проверяет следующие условия:
\begin{itemize}
  \item тренировка существует и ещё не началась;
  \item у клиента нет активной записи на эту же тренировку (запрет двойной записи);
  \item количество активных записей меньше вместимости (\texttt{capacity});
  \item новая тренировка не пересекается по времени с другими активными
        записями клиента.
\end{itemize}

Логика проверок представлена в листинге~\ref{lst:book}.
\clearpage

\begin{lstlisting}[caption={Проверка условий при записи на тренировку (workoutService.go, стр.~194--224)}, label={lst:book}]
	var existing models.Booking
	if err := s.db.Where("user_id = ? AND workout_id = ? AND status = 'active'", userID, workoutID).
		First(&existing).Error; err == nil {
		return models.Booking{}, ErrAlreadyBooked
	} else if !errors.Is(err, gorm.ErrRecordNotFound) {
		return models.Booking{}, fmt.Errorf("check existing booking: %w", err)
	}

	var workout models.Workout
	if err := s.db.Preload("Bookings", "status = 'active'").First(&workout, "id = ?", workoutID).Error; err != nil {
		if errors.Is(err, gorm.ErrRecordNotFound) {
			return models.Booking{}, ErrNotFound
		}
		return models.Booking{}, fmt.Errorf("find workout: %w", err)
	}
	if util.WorkoutStarted(workout, time.Now()) {
		return models.Booking{}, ErrWorkoutPast
	}
	if len(workout.Bookings) >= workout.Capacity {
		return models.Booking{}, ErrNoAvailableSpots
	}

	var userBookings []models.Booking
	if err := s.db.Preload("Workout").Where("user_id = ? AND status = 'active'", userID).Find(&userBookings).Error; err != nil {
		return models.Booking{}, fmt.Errorf("list user bookings: %w", err)
	}
	for _, b := range userBookings {
		if b.WorkoutID != workoutID && util.WorkoutsOverlap(workout, b.Workout) {
			return models.Booking{}, ErrTimeConflict
		}
	}
\end{lstlisting}

На клиентской стороне кнопка «Записаться» на карточке тренировки блокируется,
если мест нет, клиент уже записан или тренировка прошла. При успешной записи
отправляется запрос \texttt{POST /api/workouts/:id/book}, после чего интерфейс
обновляет состояние карточки без перезагрузки страницы. Тексты ошибок
(«Мест нет», «Вы уже записаны», «Пересечение по времени») локализованы
на русский язык в модуле \texttt{errors.js}.

\subsection{Разработка личного кабинета с историей записей и возможностью отмены}

Согласно задаче~3.5, для каждой роли пользователя реализован отдельный
личный кабинет, соответствующий макетам, разработанным в разделе~2.4.

\textbf{Кабинет клиента} (\texttt{Profile.js}) отображает сводную статистику:
общее число записей, количество предстоящих и посещённых тренировок. В разделе
«Мои записи» выводится история бронирований с указанием названия тренировки,
даты, времени и статуса. Для активных записей на будущие занятия доступна
кнопка отмены, которая вызывает эндпоинт \texttt{PUT /api/bookings/:id/cancel}
и изменяет статус записи на \texttt{cancelled}. Клиент также может редактировать
имя и загружать аватар.

\textbf{Кабинет тренера} (\texttt{TrainerProfile.js}) содержит статистику
по проведённым тренировкам и числу записавшихся клиентов, форму редактирования
профиля с полями «О себе» и «Опыт», а также загрузку фотографии. Управление
расписанием и отметка посещаемости выполняются на странице
\texttt{TrainerDashboard}.

\textbf{Кабинет администратора} (\texttt{AdminProfile.js} и
\texttt{AdminDashboard.js}) объединяет доступ к статистике платформы,
управлению пользователями и тренировками. Получение истории записей клиента
реализовано методом \texttt{MyBookings}, который загружает связанные
тренировки и данные тренера через механизм \texttt{Preload} ORM-библиотеки GORM.

\subsection{Реализация агрегации данных (статистика посещаемости, загруженность тренеров)}

В соответствии с задачей~3.6, на сервере реализован эндпоинт
\texttt{GET /api/admin/stats}, формирующий агрегированные показатели
работы фитнес-центра. Метод \texttt{GetStats} подсчитывает:
\begin{itemize}
  \item общее число активных записей (\texttt{total\_bookings});
  \item количество тренировок в системе (\texttt{total\_workouts});
  \item число тренеров (\texttt{total\_trainers}) и клиентов (\texttt{total\_clients});
  \item загруженность каждого тренера --- количество его тренировок и записей
        на них (\texttt{trainer\_load}).
\end{itemize}

Для расчёта загруженности тренеров используется SQL-запрос с группировкой
и соединением таблиц \texttt{users}, \texttt{workouts} и \texttt{bookings}
(листинг~\ref{lst:stats}). Результаты отображаются на дашборде администратора
в виде карточек с ключевыми метриками и таблицы загруженности тренеров.

\begin{lstlisting}[caption={SQL-запрос загруженности тренеров (adminService.go, стр.~159--169)}, label={lst:stats}]
	if err := s.db.Raw(`
		SELECT u.name as trainer_name,
		       COUNT(DISTINCT w.id) as workouts,
		       COUNT(DISTINCT b.id) as bookings
		FROM users u
		LEFT JOIN workouts w ON w.trainer_id = u.id
		LEFT JOIN bookings b ON b.workout_id = w.id AND b.status = 'active'
		WHERE u.role = 'trainer'
		GROUP BY u.id, u.name
		ORDER BY bookings DESC
	`).Scan(&stats.TrainerLoad).Error; err != nil {
\end{lstlisting}

Дополнительно тренер в личном кабинете видит количество своих тренировок
и уникальных записавшихся клиентов, рассчитанных на стороне клиента на основе
полученного списка занятий.

\subsection{Реализация загрузки файлов (аватар пользователя, фото тренировок)}

Задача~3.7 реализована с разделением статических и динамически загружаемых
файлов. Демонстрационные фотографии тренировок размещены в каталоге
\texttt{frontend/public/uploads/workouts/} и отдаются веб-сервером Nginx
напрямую при сборке frontend-контейнера. Файлы, загружаемые пользователями
в процессе работы (аватары и фото тренировок), сохраняются на диске серверной
части в каталоге \texttt{backend/uploads/} и доступны по URL \texttt{/uploads/...}.

Загрузка аватара выполняется запросом \texttt{POST /api/profile/avatar}.
Сервер сохраняет файл в каталог \texttt{uploads/avatars/}, генерируя уникальное
имя на основе UUID, и записывает путь в поле \texttt{avatar} таблицы
\texttt{users}. Загрузка фото тренировки доступна тренеру и администратору
через эндпоинт \texttt{POST /api/workouts/:id/photo}; файл сохраняется
в \texttt{uploads/workouts/}, а путь --- в поле \texttt{photo} таблицы
\texttt{workouts}.

Вспомогательная функция сохранения файла представлена в листинге~\ref{lst:upload}.

\begin{lstlisting}[caption={Сохранение загруженного файла (upload.go, стр.~11--17)}, label={lst:upload}]
func saveUploadedFile(file *multipart.FileHeader, dir string, c *gin.Context) (string, error) {
	filename := uuid.New().String() + filepath.Ext(file.Filename)
	if err := c.SaveUploadedFile(file, dir+"/"+filename); err != nil {
		return "", err
	}
	return filename, nil
}
\end{lstlisting}

Для раздачи файлов backend регистрирует статический маршрут
\texttt{r.Static("/uploads", "./uploads")}. В production-окружении Nginx
сначала ищет запрошенный файл среди статических ресурсов frontend, а при
отсутствии --- проксирует запрос на backend, что обеспечивает корректную
работу как демонстрационных, так и пользовательских изображений. Каталог
\texttt{backend/uploads} монтируется как Docker-volume, поэтому загруженные
файлы сохраняются при перезапуске контейнеров.

Таким образом, в третьей главе описана практическая реализация всех задач
раздела~3 плана работ: от базовой структуры интерфейса до загрузки файлов
и формирования статистики. Разработанное приложение FitCenter развёртывается
с помощью \textbf{Docker Compose} и включает три сервиса: frontend (Nginx +
React), backend (Go + Gin) и СУБД PostgreSQL.
